\chapter{Khảo sát hệ thống thông tin}
\section{Mục tiêu và cách tiến hành}
Bảng hỏi tập trung vào HTTT hiện tại: công cụ, dữ liệu, quy trình đơn hàng/tồn kho/khách hàng/nhân sự, phân quyền, báo cáo và khó khăn. Đây là bộ câu hỏi đề xuất, chưa có câu trả lời thực tế. Nhóm cần chọn người trả lời có hiểu biết về quy trình và ghi nhận mẫu khảo sát trước khi viết kết luận.

\section{Phiếu khảo sát đề xuất}
Với câu chọn nhiều phương án, ghi rõ tất cả phương án được chọn. Không điền thông tin cá nhân nếu không cần thiết.
\begin{enumerate}
\item Hiện doanh nghiệp sử dụng những phần mềm hoặc công cụ nào để quản lý hoạt động?
\item Dữ liệu khách hàng hiện được lưu ở đâu và ai được quyền cập nhật?
\item Doanh nghiệp có quy định mã định danh khách hàng thống nhất không? Nếu có, mã được tạo ở bước nào?
\item Quy trình tiếp nhận một đơn hàng từ lúc khách đặt đến khi hoàn tất gồm những bước nào?
\item Đơn hàng phát sinh từ những kênh nào và dữ liệu giữa các kênh được đồng bộ ra sao?
\item Nhân viên kiểm tra tình trạng thanh toán bằng cách nào? Ai được phép xác nhận đã thanh toán?
\item Quy trình cập nhật tồn kho khi nhập hàng, bán hàng, đổi trả hoặc kiểm kê hiện nay như thế nào?
\item Làm thế nào để tìm nhanh sản phẩm theo tên, mã, danh mục, size, màu hoặc nhà cung cấp?
\item Doanh nghiệp theo dõi thông tin và lịch sử tương tác khách hàng ở những hệ thống nào?
\item Khách hàng có thể yêu cầu cập nhật hoặc xóa thông tin cá nhân bằng cách nào?
\item Phản hồi/khiếu nại của khách được tiếp nhận, phân loại, giao xử lý và đóng như thế nào?
\item Doanh nghiệp đang gửi khảo sát khách hàng bằng kênh nào? Có theo dõi người nhận và trạng thái trả lời không?
\item Những báo cáo về doanh thu, đơn hàng, khách hàng, tồn kho hoặc nhân sự nào được sử dụng thường xuyên?
\item Báo cáo hiện được tổng hợp thủ công hay tự động? Mất bao lâu để có báo cáo cần thiết?
\item Tài khoản người dùng được cấp, khóa và thu hồi quyền theo quy trình nào?
\item Dữ liệu được sao lưu với tần suất nào? Đã thử phục hồi bản sao lưu gần đây chưa?
\item Khó khăn hoặc lỗi nào thường gặp khi tìm kiếm, nhập liệu, đồng bộ hoặc lập báo cáo?
\item Theo người trả lời, thay đổi nào trong HTTT sẽ giúp quy trình hiện tại tốt hơn? Vì sao?
\end{enumerate}

\section{Mẫu tổng hợp kết quả}
\begin{longtable}{|p{0.07\textwidth}|p{0.42\textwidth}|p{0.18\textwidth}|p{0.22\textwidth}|}\hline
\textbf{Câu}&\textbf{Chủ đề}&\textbf{Số trả lời hợp lệ}&\textbf{Kết quả/tỷ lệ thật}\\\hline
1--3&Công cụ và quản lý dữ liệu&[CẦN BỔ SUNG]&[CẦN BỔ SUNG]\\\hline
4--8&Đơn hàng và tồn kho&[CẦN BỔ SUNG]&[CẦN BỔ SUNG]\\\hline
9--12&Khách hàng, phản hồi và khảo sát&[CẦN BỔ SUNG]&[CẦN BỔ SUNG]\\\hline
13--14&Báo cáo&[CẦN BỔ SUNG]&[CẦN BỔ SUNG]\\\hline
15--16&Bảo mật và sao lưu&[CẦN BỔ SUNG]&[CẦN BỔ SUNG]\\\hline
17--18&Khó khăn và nhu cầu cải tiến&[CẦN BỔ SUNG]&[CẦN BỔ SUNG]\\\hline
\end{longtable}
Tổng số người tham gia: [CẦN BỔ SUNG]. Số phiếu hợp lệ: [CẦN BỔ SUNG]. Thời gian/phương thức khảo sát: [CẦN BỔ SUNG]. Tỷ lệ câu trả lời = số lượt chọn phương án / số phiếu hợp lệ có trả lời câu hỏi; nếu câu hỏi cho phép chọn nhiều, tổng tỷ lệ có thể vượt 100\%.

\section{Kết luận khảo sát}
[CẦN BỔ SUNG: chỉ viết kết luận sau khi thu thập dữ liệu thật; nêu điểm mạnh, hạn chế được nhiều người xác nhận, nhu cầu ưu tiên và giới hạn mẫu khảo sát. Không dùng số liệu seed/demo để thay thế.]
